\section{Conclusion}

We present AutoQ-VIS, a quality-aware self-training framework that advances unsupervised video instance segmentation through iterative pseudo-label refinement with automatic quality control. By establishing a closed-loop system of pseudo-label generation and automatic quality assessment, our method achieves state-of-the-art performance (52.6 $\text{AP}_{50}$) on YouTubeVIS-2019 \texttt{val} split without requiring any human annotations. The simple quality predictor proves effective in pseudo-label quality assessment. 
% Future work may improve pseudo-label selection via temporal-aware quality scoring that leverages inter-frame consistency.    

% \onur{I did not like the way you formulate the main contribution: "bridging the synthetic-to-real domain gap". This sounds like a domain adaptation paper but our method is about sth. else. Can we rephrase it?}


% Current limitations stem from two aspects: (1) Hyperparameter sensitivity—the fixed quality threshold $\tau_{th}$ and the number of self-training rounds require careful tuning across datasets; (2) Simplistic quality modeling—our predictor currently evaluates frame-level mask quality without considering temporal consistency. Future work may explore temporal-aware quality scoring and online threshold adaptation to further improve performance.


% \onur{refine the references, inconsistent formating}